\section{Introduction}
\label{sec:intro}

The ear is a useful biometric trait in surveillance and forensic settings because it can be photographed at a distance and without the cooperation of the subject~\cite{emersic2017ear}. In such footage the ear often occupies a few dozen pixels. In EarVN1.0~\cite{hoang2019earvn}, a collection of ear crops gathered from the web, several thousand images have a short side of 24 to 48 pixels, and all of them are stored with lossy compression. Low resolution and compression are both known to affect ear recognition~\cite{emersic2019uerc,rathgeb2016compression}. An operator who inspects such an image, or a pipeline that requires a larger input, has to enlarge it first, and in interactive or video applications this has to be done in real time.

Efficient single-image super-resolution (SR) is a natural candidate for this step. Work on efficient SR, in recent years organised around the NTIRE Efficient SR challenges~\cite{ren2024ntire,ren2025ntire,ren2026ntire}, has produced networks with a few hundred thousand parameters that run in milliseconds. These networks are designed and ranked under a single protocol: large natural images, bicubic downsampling by a factor of four, and no compression. JPEG-compressed inputs of a few dozen pixels differ from this protocol in content, size and degradation. The one study of SR for ear images that we found~\cite{markicevic2023ear} trains two general-purpose networks and evaluates recognition; it does not address efficient models. It is therefore unclear whether a published efficient SR model is useful on small ear images, and, if it is not, whether the architecture or the training data should be changed.

It is known that SR models trained with bicubic downsampling lose accuracy on compressed inputs~\cite{lugmayr2019unsupervised,chen2018cisrdcnn}, and training with compression is an established remedy~\cite{zhang2021bsrgan,wang2021realesrgan}. This paper does not revisit that general result. It addresses three questions that, to our knowledge, have not been examined for small images of a biometric trait. The first is at which compression level published efficient models stop improving on interpolation, and what determines this level. The second is whether the effect can be observed on real small images, for which no ground truth exists. The third is what is required to adapt a published real-time model to such inputs.

For the first question, we build a benchmark of ear images with clean ground truth at three input sizes from the AMI database, add two in-the-wild datasets, EarVN1.0 and AWEx, and evaluate \NPub{} fidelity-oriented SR networks with public weights, with a control on natural images. In each of the \RatioCells{} settings we tested (four datasets, inputs of \RatioPxMin{} to \RatioPxMax{} pixels, JPEG quality 60 to 93), each model reduces the interpolation error (to \ReduceMin{} to \ReduceMax{} of its value) and amplifies the error added by compression (by a factor of \AmpMin{} to \AmpMax{}). The models fall below bicubic interpolation when the second error exceeds a fraction of the first that lies between \CrossAllMin{} and \CrossAllMax{}, depending on the dataset, and input size has little influence on this point (Section~\ref{sec:natural}). Ear images are smooth and their interpolation error is small. Consistent with this, the gain changes sign at JPEG quality \CrossQAmiMin{} to \CrossQAmiMax{} on AMI, compared with \CrossQNatMin{} to \CrossQNatMax{} on natural images. The two quality levels that account for \QLowShare{}\% and \QHighShare{}\% of the real small ear images we examined, 75 and 93, lie on opposite sides of this point: at quality 75 all \NPub{} models are \AmiQLowGapMin{} to \AmiQLowGapMax{}~dB below bicubic interpolation on AMI, and at quality 93 all of them remain above it (Section~\ref{sec:analysis}).

For the second question, we use two properties of real image files that do not require ground truth: the JPEG quality stored in each file, and the identity label. On the real small images of EarVN1.0, the published SPAN weights change the rank-1 identification rate by \RecQLowPub{} percentage points relative to bicubic upscaling for probes stored at quality 70 to 80, and by \RecQHighPub{} points for probes stored at quality 90 or above, which is consistent with the results on synthesised inputs. A model trained with compression changes it by \RecQLowEst{} and \RecQHighEst{} points. Over all probes the rate is \RecRankBic{}\% with bicubic upscaling, \RecRankPub{}\% with the published weights and \RecRankEst{}\% with the model trained with compression, and a second recogniser gives the same ordering. On AWEx no upscaling method changes the rate significantly. AWEx differs from EarVN1.0 in two respects: its small images show no JPEG block structure, and the recognisers lose much less accuracy between large and small images. We examined these differences after obtaining the result, and the two datasets do not allow them to be separated. We report them as possible conditions for an improvement (Section~\ref{sec:recog}).

For the third question, we compare the effect of the architecture with that of the training degradation. Under JPEG compression at quality 75 the \NPub{} architectures span \SpreadQLow{}~dB, compared with \SpreadClean{}~dB on clean inputs, and a network with \ParamsRatio{} times more parameters does not perform better than a \ParamsSpan{}M-parameter network. Changing the degradation used to fine-tune one backbone changes its PSNR on such inputs by \DegEffMin{} to \DegEffMax{}~dB. We therefore keep a published real-time backbone and change its training. We measure the degradation of real small ear images (the JPEG quality from the quantisation tables stored in the files, the noise level with a standard estimator, and a mild blur selected to match sharpness statistics), and an ablation indicates that the range of JPEG quality levels is the component that matters most. We also report that a model fine-tuned on laboratory images alone fails on bright in-the-wild images (\CollapsePsnr{}~dB, compared with \CollapsePub{}~dB for the weights it was initialised from), and describe two measures that prevented this in our experiments (Section~\ref{sec:stable}).

\paragraph{Contributions}
\begin{enumerate}
\item An evaluation of \NPub{} published efficient SR models on a benchmark of small ear images (clean ground truth at three input sizes, subject-wise cross-validation, two in-the-wild datasets), with a control on natural images. The evaluation relates the advantage of these models over interpolation to the ratio of compression error to interpolation error, and shows that the gain changes sign at a lighter compression for ear images than for natural images.
\item An evaluation on real small ear images without ground truth, based on the JPEG quality stored in the image files and on identity labels, with two datasets and two recognisers. On EarVN1.0 the results agree with those on synthesised inputs. On AWEx no upscaling method changes identification, and we report two differences between the datasets as possible explanations.
\item A training procedure for adapting a published real-time model without architectural changes: a degradation measured from real small ear images, an ablation of its components, and two measures that kept fine-tuning stable on in-the-wild data. The resulting model improves on bicubic interpolation by \EstVsBicMin{} to \EstVsBicMax{}~dB and on the published weights by \EstVsPubMin{} to \EstVsPubMax{}~dB on three datasets.
\end{enumerate}

We do not propose a new architecture. The results have several limits, which we discuss in Section~\ref{sec:discussion}. The model is \EstVsPubCleanMin{} to \EstVsPubCleanMax{}~dB below the published weights on uncompressed inputs. Large perceptual models have a lower LPIPS and a statistically similar identification rate, with a fidelity below bicubic interpolation and \LatRatio{} times the latency on a desktop GPU. The gain is associated with covering the range of compression levels: variants trained with JPEG compression alone give results on real images that we could not distinguish from those of the full measured degradation (Sections~\ref{sec:ablation} and~\ref{sec:recog}).
